\section{总结与延伸}

\subsection{全讲总结表}
\begin{table}[H]
\centering
\renewcommand{\arraystretch}{1.2}
\begin{tabular}{p{2.7cm}p{3.9cm}p{4.2cm}}
\toprule
\textbf{阶段} & \textbf{核心问题} & \textbf{本讲给出的可执行答案} \\
\midrule
Pre-training & 底座能力如何在预算内做强 & 用可复现 OLMo 配方优化数据与训练细节 \\
Post-training & 如何把能力变成可用行为与推理优势 & SFT 打协议，DPO/PPO 调偏好，RLVR 冲可验证推理 \\
Inference-time & 如何在部署端进一步提效 & 按请求价值使用 test-time scaling 与检索校验闭环 \\
开放生态 & 如何让改进可累积 & 开放数据/代码/评估链路，强化可复现与可审计 \\
\bottomrule
\end{tabular}
\caption{Open Training Recipes for Reasoning 的结构化结论}
\end{table}

\subsection{核心结论}
\begin{importantbox}{一句话总结}
Reasoning 能力是系统工程结果，不是单点技巧；最有效路径是 ``可复现预训练 + 分阶段后训练 + 预算化测试时扩展'' 的组合配方。
\end{importantbox}

\begin{knowledgebox}{本讲对实践者最有价值的三点}
\begin{enumerate}
    \item 先构建可复现实验链路，再追求更高分数；
    \item 用可验证奖励扩大 RL 训练可自动化范围；
    \item 将 test-time scaling 作为部署预算管理能力，而不是临时补丁。
\end{enumerate}
\end{knowledgebox}

\subsection{进一步阅读}
\begin{itemize}
    \item OLMo / OLMo 2 技术报告与训练代码（AI2）
    \item Tulu 3 后训练配方论文与开源实现
    \item DPO / PPO 系列偏好优化论文与长度偏置修正方法
    \item RLVR 与 DeepSeek-R1 路线相关公开资料
    \item Self-RAG 与 test-time scaling 相关研究
\end{itemize}

\subsection{延伸问题}
\begin{warningbox}{值得继续追踪的开放问题}
\begin{itemize}
    \item 如何为非数学非代码任务构建低成本可信验证器？
    \item 如何建立跨模型、跨数据版本的一致评估协议？
    \item 如何在固定预算下动态决定 ``再采样'' 是否值得？
\end{itemize}
\end{warningbox}

%% --- Fin del contenido --- %%

\end{document}
